% Standalone, submission-ready write-up of the two empirical questions.
% Compile from hw_empirical/ with: pdflatex empirical_answers.tex
% The two answer bodies can also be copied into the corresponding \ans locations.
\documentclass[10pt]{article}
\usepackage[letterpaper,margin=0.67in]{geometry}
\usepackage[T1]{fontenc}
\usepackage{amsmath}
\usepackage{booktabs}
\usepackage{graphicx}
\usepackage{hyperref}
\usepackage{microtype}
\hypersetup{colorlinks=true,urlcolor=blue}
\setlength{\parindent}{0pt}
\setlength{\parskip}{0.45em}

\begin{document}

\section*{Empirical: Dense vs. Mixture-of-Experts Serving}

I benchmarked the three checkpoints with vLLM 0.25.0 and PyTorch 2.11.0 on one
NVIDIA H100 80GB.  All runs used tensor parallelism 1, online FP8 weight
quantization, an FP8 KV cache with dynamically calculated scales, and disabled
prefix caching.  Tokenization was excluded from timing.  For the prefill
workload, I generated one token from a single input and varied its length.  For
the decode workload, I used 32-token prompts, forced each response to contain
512 new tokens, and varied the number of simultaneous responses.  Each point
below is the mean of three timed trials after one warmup; the error bars are one
sample standard deviation.

\begin{figure}[ht]
  \centering
  \includegraphics[width=0.72\linewidth]{artifacts/serving/figures/prefill_throughput.pdf}
  \caption{Prefill throughput as input context length changes.}
  \label{fig:prefill}
\end{figure}

\begin{table}[ht]
\centering
\small
\begin{tabular}{rrrr}
\toprule
Input tokens & OLMo-2 1B dense & OLMoE 1B/7B & OLMo-2 7B dense \\
\midrule
128  & $31.9\!\pm\!1.3$ & $29.1\!\pm\!0.9$ & $17.5\!\pm\!0.5$ \\
512  & $104.4\!\pm\!0.7$ & $74.7\!\pm\!0.2$ & $46.1\!\pm\!0.1$ \\
2048 & $204.2\!\pm\!0.6$ & $125.7\!\pm\!0.4$ & $56.5\!\pm\!0.2$ \\
3072 & $217.2\!\pm\!0.5$ & $164.0\!\pm\!0.2$ & $54.8\!\pm\!0.1$ \\
\bottomrule
\end{tabular}
\caption{Mean prefill throughput in thousands of input tokens/s ($\pm$ one SD).}
\end{table}

\begin{figure}[ht]
  \centering
  \includegraphics[width=0.72\linewidth]{artifacts/serving/figures/decode_throughput.pdf}
  \caption{Aggregate decode throughput as concurrency changes.}
  \label{fig:decode}
\end{figure}

\begin{table}[ht]
\centering
\small
\begin{tabular}{rrrr}
\toprule
Parallel generations & OLMo-2 1B dense & OLMoE 1B/7B & OLMo-2 7B dense \\
\midrule
1  & $750\!\pm\!0$       & $621\!\pm\!0$       & $244\!\pm\!0$ \\
2  & $1470\!\pm\!0$      & $1204\!\pm\!0$      & $485\!\pm\!0$ \\
4  & $2914\!\pm\!1$      & $2238\!\pm\!0$      & $955\!\pm\!1$ \\
8  & $5718\!\pm\!13$     & $4377\!\pm\!9$      & $1799\!\pm\!0$ \\
16 & $10493\!\pm\!2$     & $8061\!\pm\!21$     & $3285\!\pm\!2$ \\
32 & $17436\!\pm\!22$    & $13807\!\pm\!3$     & $5312\!\pm\!3$ \\
\bottomrule
\end{tabular}
\caption{Mean decode throughput in output tokens/s ($\pm$ one SD), rounded to
the nearest token/s; a displayed SD of zero is smaller than 0.5 token/s.}
\end{table}

Across every tested operating point, OLMoE was between the two dense models:
there was no crossover at which it beat the 1B dense model, but it always beat
the 7B dense model.  In prefill, OLMoE was 9--38\% slower than the 1B dense
model but 62--199\% faster than the 7B dense model.  In decode, it was 17--23\%
slower than the 1B model and 134--160\% faster than the 7B model.  At 32
parallel generations, for example, throughput was 17,436, 13,807, and 5,312
tokens/s for the small dense, MoE, and large dense models respectively.

Thus, OLMoE behaves more like the small model in arithmetic cost: routing only
about 1B active parameters per token makes its throughput substantially closer
to the 1B dense checkpoint than to the 7B dense checkpoint.  It retains a
large-model-like cost because its roughly 7B stored expert parameters must fit
in memory and the system must route tokens and access different expert weights;
this is consistent with its persistent throughput deficit relative to the 1B
dense model.  Increasing concurrency improved aggregate decode throughput for
all models by batching more work, without changing their ordering.  Prefill
throughput initially increased as longer inputs amortized fixed launch and
scheduling overheads; the 7B dense model saturated near 2048 tokens while the
two smaller-active-compute models continued to improve through the tested
range.

The implementation is
\path{hw_empirical/serving/benchmark_vllm.py}.  All individual repetitions and
run metadata are in \path{hw_empirical/artifacts/serving/raw/}; the plotted
means are in
\path{hw_empirical/artifacts/serving/figures/serving_summary.csv}.

\clearpage
\section*{Empirical: Supervised Fine-Tuning with LoRA}

I trained \path{Qwen/Qwen2.5-0.5B} for one epoch on a fixed set of 1,000
\path{train_sft} conversations and evaluated on the first 100
\path{test_sft} conversations.  Every run used the same H100 GPU, examples,
seed, sequence length 512, effective batch size 16, learning rate $5\times
10^{-5}$, BF16 precision, and evaluation schedule.  LoRA used PEFT's
\path{target_modules="all-linear"}; full tuning updated all weights.

\begin{table}[ht]
\centering
\small
\begin{tabular}{lrrrrr}
\toprule
Variant & Trainable params. & Peak GPU GiB & Time (s) & Mean train loss & Final val. loss \\
\midrule
LoRA $r=1$  & 549,888     & 3.703 & 79.40 & 1.954 & 1.816 \\
LoRA $r=4$  & 2,199,552   & 3.732 & 77.23 & 1.874 & 1.751 \\
LoRA $r=16$ & 8,798,208   & 3.857 & 78.22 & 1.816 & 1.730 \\
Full         & 494,032,768 & 6.486 & 39.72 & 1.748 & 1.665 \\
\bottomrule
\end{tabular}
\caption{SFT results. Peak GPU memory is PyTorch peak reserved memory, and time is
the trainer's training runtime.}
\label{tab:sft}
\end{table}

\begin{figure}[ht]
  \centering
  \includegraphics[width=0.92\linewidth]{artifacts/sft/report/sft_loss_curves.pdf}
  \caption{Training and held-out validation losses over the 63 optimizer steps.}
  \label{fig:sft-loss}
\end{figure}

The LoRA parameter counts scaled linearly with rank and matched the analytical
counts.  Peak reserved memory changed only from 3.703 GiB at rank 1 to 3.857
GiB at rank 16 because frozen base weights and activations dominate the small
adapter allocations.  The LoRA runs used 40--43\% less peak memory than full
tuning.  In this particular H100 implementation, however, they took about
1.94--2.00 times as long as full tuning.  LoRA therefore reduced memory but not
wall-clock time here; this likely reflects adapter/wrapper overhead and kernel
efficiency while the small full model fit comfortably.  This timing
is an implementation-specific measurement, not a claim that full tuning is
generally faster.

The loss curves show a consistent capacity ordering.  From step 10 to the end,
validation loss fell from 2.001 to 1.816 for rank 1, 1.940 to 1.751 for rank 4,
1.833 to 1.730 for rank 16, and 1.754 to 1.665 for full tuning.  Training loss
was noisier across logging windows but had the same overall downward trend and
ordering.  Validation loss continued decreasing or remained nearly flat at the
end, so there was no clear evidence of overfitting during this single epoch.
Higher LoRA rank helped, although the improvement from rank 4 to 16 was smaller
than that from rank 1 to 4, and full tuning achieved the lowest losses.

For a qualitative comparison I greedily generated at most 192 tokens from the
base model and all four tuned variants on the same five held-out user prompts.
Fine-tuning produced a meaningful but inconsistent change:
\begin{itemize}
  \item On the free-verse prompt, the base model explained resilience in prose.
  The tuned models instead began with concrete imagery, and the full model used
  line-broken verse centered on climbing a mountain and reaching its summit.
  \item On the affordable-dental-care prompt, the base model opened as a letter
  (``Dear [Recipient]'').  Rank 16 and full tuning opened as speeches (``Good
  morning, everyone''), following the requested format more directly.  Rank 1
  still resembled a letter, while rank 4 adopted a speech opening but became
  repetitive.
  \item The improvement was not uniform.  Several variants repeated phrases or
  hallucinated details; some science-answer completions emitted stray role or
  template-like text, and all historical-fiction variants contained questionable
  historical details.  One epoch on 1,000 examples therefore improved visible
  instruction-following and conversational formatting, especially at higher
  rank/full tuning, but did not make this 0.5B base model reliably correct or
  fluent.
\end{itemize}

The implementation is \path{hw_empirical/sft/train_sft.py}.  Raw per-step
losses and metrics are under \path{hw_empirical/artifacts/sft/runs/}; the shared
data fingerprints are in \path{hw_empirical/artifacts/sft/data/manifest.json},
and all matched generations are in
\path{hw_empirical/artifacts/sft/report/qualitative_generations.jsonl}.

\end{document}
